\section{Conclusions}

La planificació de tasques asíncrones distribueix una capacitat limitada. FIFO
ofereix una línia base simple; la prioritat estricta pot accelerar treballs
urgents però exposa les classes baixes a esperes extremes; i l'envelliment
intenta moderar aquest risc a canvi d'alterar el servei de les prioritats
inicials. La literatura no permet declarar una política universalment superior,
perquè els resultats depenen de la càrrega, les durades, les classes i la
mètrica observada.

La simulació d'esdeveniments discrets permet comparar alternatives sense
dependre del temps real ni d'una infraestructura productiva. Totes les
polítiques han de rebre entrades equivalents, els experiments aleatoris han de
conservar llavors i repeticions, i les mètriques s'han d'analitzar globalment i
per classes. La verificació del motor, la delimitació de la validesa i la
documentació dels experiments formen part de la contribució.

Les eines existents cobreixen fragments del problema, però no s'ha identificat
una solució directament equivalent a la plataforma proposada. Aquest resultat
justifica un prototip propi de menor abast, amb polítiques intercanviables,
escenaris sintètics persistits i una visualització orientada a interpretar els
resultats. L'estat de l'art estableix així la base conceptual i metodològica per
prendre les decisions de disseny i avaluar-les sense anticipar conclusions.
